% Options for packages loaded elsewhere
\PassOptionsToPackage{unicode}{hyperref}
\PassOptionsToPackage{hyphens}{url}
\documentclass[
]{article}
\usepackage{xcolor}
\usepackage[margin=1in]{geometry}
\usepackage{amsmath,amssymb}
\setcounter{secnumdepth}{-\maxdimen} % remove section numbering
\usepackage{iftex}
\ifPDFTeX
  \usepackage[T1]{fontenc}
  \usepackage[utf8]{inputenc}
  \usepackage{textcomp} % provide euro and other symbols
\else % if luatex or xetex
  \usepackage{unicode-math} % this also loads fontspec
  \defaultfontfeatures{Scale=MatchLowercase}
  \defaultfontfeatures[\rmfamily]{Ligatures=TeX,Scale=1}
\fi
\usepackage{lmodern}
\ifPDFTeX\else
  % xetex/luatex font selection
\fi
% Use upquote if available, for straight quotes in verbatim environments
\IfFileExists{upquote.sty}{\usepackage{upquote}}{}
\IfFileExists{microtype.sty}{% use microtype if available
  \usepackage[]{microtype}
  \UseMicrotypeSet[protrusion]{basicmath} % disable protrusion for tt fonts
}{}
\makeatletter
\@ifundefined{KOMAClassName}{% if non-KOMA class
  \IfFileExists{parskip.sty}{%
    \usepackage{parskip}
  }{% else
    \setlength{\parindent}{0pt}
    \setlength{\parskip}{6pt plus 2pt minus 1pt}}
}{% if KOMA class
  \KOMAoptions{parskip=half}}
\makeatother
\usepackage{graphicx}
\makeatletter
\newsavebox\pandoc@box
\newcommand*\pandocbounded[1]{% scales image to fit in text height/width
  \sbox\pandoc@box{#1}%
  \Gscale@div\@tempa{\textheight}{\dimexpr\ht\pandoc@box+\dp\pandoc@box\relax}%
  \Gscale@div\@tempb{\linewidth}{\wd\pandoc@box}%
  \ifdim\@tempb\p@<\@tempa\p@\let\@tempa\@tempb\fi% select the smaller of both
  \ifdim\@tempa\p@<\p@\scalebox{\@tempa}{\usebox\pandoc@box}%
  \else\usebox{\pandoc@box}%
  \fi%
}
% Set default figure placement to htbp
\def\fps@figure{htbp}
\makeatother
\setlength{\emergencystretch}{3em} % prevent overfull lines
\providecommand{\tightlist}{%
  \setlength{\itemsep}{0pt}\setlength{\parskip}{0pt}}
<style>
body { font-family: sans-serif; line-height: 1.5; }
h1, h2, h3 { margin-top: 1.4em; }
table { width: auto; }
</style>
\usepackage{bookmark}
\IfFileExists{xurl.sty}{\usepackage{xurl}}{} % add URL line breaks if available
\urlstyle{same}
\hypersetup{
  pdftitle={Problem Set 1: Conditioning on Observables and Matching},
  pdfauthor={Advanced Observational Causal Inference},
  hidelinks,
  pdfcreator={LaTeX via pandoc}}

\title{Problem Set 1: Conditioning on Observables and Matching}
\author{Advanced Observational Causal Inference}
\date{}

\begin{document}
\maketitle

\section{Problem Set 1: Conditioning on Observables and
Matching}\label{problem-set-1-conditioning-on-observables-and-matching}

\subsection{Purpose}\label{purpose}

This problem set asks you to evaluate whether conditioning on observed
covariates can support a causal interpretation in an observational
study. You will work with Arceneaux, Gerber, and Green (2010), ``A
Cautionary Note on the Use of Matching to Estimate Causal Effects: An
Empirical Example Comparing Matching Estimates to an Experimental
Benchmark.'' The paper uses a large-scale voter-mobilization experiment
to compare an observational matching estimate with an experimental
benchmark.

The central methodological question is:

\begin{quote}
\textbf{When does conditioning on observed covariates allow us to
recover a causal effect, and what can empirical diagnostics tell us when
the identifying assumption is fundamentally untestable?}
\end{quote}

You should approach the paper first as an observational researcher. In
the middle of the assignment, you will then use the randomized component
of the study as an experimental benchmark. Do not use that benchmark to
justify your earlier answers retrospectively; the point is to assess
what could have been learned from the observational analysis alone.

\subsection{Data and replication
materials}\label{data-and-replication-materials}

Use the replication materials supplied with the assignment. All
empirical work should be conducted in \textbf{R} and submitted as a
reproducible R Markdown document or equivalent Quarto document.

The paper and replication materials are:

\begin{itemize}
\tightlist
\item
  Arceneaux, Kevin, Alan S. Gerber, and Donald P. Green. 2010. ``A
  Cautionary Note on the Use of Matching to Estimate Causal Effects: An
  Empirical Example Comparing Matching Estimates to an Experimental
  Benchmark.'' \emph{Sociological Methods \& Research}.
\end{itemize}

You may consult the paper while completing the assignment. You may also
consult the course readings on causal identification, DAGs, matching,
and weighting.

\subsection{General submission
requirements}\label{general-submission-requirements}

Your submission should contain publication-quality tables and figures.
Do not paste raw console output into the assignment. Tables and figures
should have informative labels, appropriate numbers of digits, sample
sizes where relevant, and captions that explain what is being shown.

For each empirical result, explain what the quantity represents and why
the analysis is relevant to the causal question. A correct numerical
result without an explanation of its causal interpretation will receive
limited credit.

You may use AI tools for coding assistance, debugging, and consultation,
but you remain responsible for every line of code, statistical result,
and causal claim in your submission. Be prepared to explain your
analysis without relying on the AI system that helped produce it.

\begin{center}\rule{0.5\linewidth}{0.5pt}\end{center}

\section{Part I. Formalize the causal question and identification
strategy}\label{part-i.-formalize-the-causal-question-and-identification-strategy}

\subsubsection{1. Define the causal estimand. {[}5
points{]}}\label{define-the-causal-estimand.-5-points}

AGG distinguish between treatment assignment, actual receipt of the
treatment, and whether a subject is reachable by phone. Using the
notation in the paper, define precisely:

\begin{itemize}
\tightlist
\item
  the treatment of substantive interest;
\item
  the outcome;
\item
  the target population for the treatment effect; and
\item
  the causal estimand the authors call the treatment-on-treated effect
  (ATT).
\end{itemize}

Explain why, in this design, the ATT can also be interpreted as the
average treatment effect among compliers.

\subsubsection{2. Reconstruct the observational identification argument.
{[}10
points{]}}\label{reconstruct-the-observational-identification-argument.-10-points}

Suppose you did \textbf{not} know that the data came from a randomized
experiment. You simply observed whether a person was contacted by the
phone bank, their 2004 turnout, and the pre-treatment covariates used
for matching.

Construct a DAG representing the causal argument required for a matching
analysis to identify the ATT.

Your DAG should include, at minimum:

\begin{itemize}
\tightlist
\item
  treatment/contact;
\item
  2004 turnout;
\item
  observed pre-treatment covariates used in the matching procedure; and
\item
  at least one unobserved characteristic that could plausibly affect
  both treatment receipt and turnout.
\end{itemize}

Then explain in words what causal pathways the matching strategy is
attempting to block.

\subsubsection{3. State the identifying assumption. {[}10
points{]}}\label{state-the-identifying-assumption.-10-points}

Write down a formal version of the conditional ignorability assumption
required for the observational comparison to identify the ATT.

Then answer:

\begin{enumerate}
\def\labelenumi{\arabic{enumi}.}
\tightlist
\item
  What does this assumption imply about the untreated potential outcome
  among people who were reachable versus those who were not?
\item
  Why does exact matching on observed covariates not establish that this
  assumption is true?
\item
  Using your DAG from Question 2, explain what happens to identification
  if an unobserved cause of both reachability and turnout remains.
\end{enumerate}

\subsubsection{4. Backdoor criterion. {[}10
points{]}}\label{backdoor-criterion.-10-points}

Using the DAG from Question 2, assess whether the proposed adjustment
set satisfies the backdoor criterion.

Identify which paths are blocked by conditioning on the observed
covariates and which paths would remain open if an important confounder
were unobserved.

Your answer should distinguish between:

\begin{itemize}
\tightlist
\item
  a problem with the \textbf{choice of estimator}; and
\item
  a problem with the \textbf{identification assumption itself}.
\end{itemize}

\begin{center}\rule{0.5\linewidth}{0.5pt}\end{center}

\section{Part II. Reproduce the observational
analysis}\label{part-ii.-reproduce-the-observational-analysis}

In the following questions, put yourself in the position of a researcher
who does not have access to the experimental benchmark.

\subsubsection{5. Reproduce the matching sample. {[}8
points{]}}\label{reproduce-the-matching-sample.-8-points}

Using the covariates described in the paper, reconstruct the
exact-matching procedure used for the GOTV treatment and the matched
control group.

Report:

\begin{itemize}
\tightlist
\item
  the number of treated observations retained;
\item
  the number of matched controls retained;
\item
  the number of exact-matching strata; and
\item
  a balance table comparing the treated and matched control groups.
\end{itemize}

Your balance table should resemble the presentation in Table 2 of the
paper.

Briefly explain what the exact-match design guarantees about the
observed covariates.

\subsubsection{6. Reproduce the observational estimate. {[}10
points{]}}\label{reproduce-the-observational-estimate.-10-points}

Reproduce the observational matching estimate for the effect of the GOTV
phone call on 2004 turnout.

Report the point estimate and its standard error, and compare your
result to the estimate reported in the paper.

Interpret the coefficient in percentage-point terms.

Then repeat the analysis for the buckle-up placebo treatment.

\subsubsection{7. What does balance establish? {[}7
points{]}}\label{what-does-balance-establish-7-points}

The authors emphasize that their matched treatment and control groups
are perfectly balanced on the observed covariates.

Explain precisely what this evidence establishes and what it does
\textbf{not} establish.

Your answer should address the distinction between:

\begin{itemize}
\tightlist
\item
  balance on observed covariates;
\item
  conditional ignorability; and
\item
  causal identification.
\end{itemize}

\begin{center}\rule{0.5\linewidth}{0.5pt}\end{center}

\section{Part III. Diagnose the observational
design}\label{part-iii.-diagnose-the-observational-design}

\subsubsection{8. Placebo treatment test. {[}10
points{]}}\label{placebo-treatment-test.-10-points}

The paper includes a second treatment condition: a phone call
encouraging people to buckle their seat belts.

Use the replication data to reproduce the authors' placebo-treatment
analysis.

Then explain:

\begin{enumerate}
\def\labelenumi{\arabic{enumi}.}
\tightlist
\item
  Why is the buckle-up message a useful placebo treatment for voter
  turnout?
\item
  What identifying assumption would have to hold for a correctly
  implemented matching analysis to produce a zero effect here?
\item
  What does the large positive matching estimate imply about the
  credibility of the observational design?
\item
  Does a successful placebo test prove conditional ignorability? Explain
  why or why not.
\end{enumerate}

\subsubsection{9. Lagged outcome placebo test. {[}10
points{]}}\label{lagged-outcome-placebo-test.-10-points}

Following the paper's diagnostic exercise, use a pre-treatment turnout
measure as the outcome. Reconstruct the authors' matching strategy using
the reduced set of matching covariates described in the paper, then
estimate the relationship between treatment/contact and the earlier
voting outcome.

Report the result for the GOTV treatment and explain why this is a
useful diagnostic.

Why can a correlation between treatment receipt and a pre-treatment
outcome not be interpreted as a causal effect of the phone call?

What does the result suggest about selection into treatment?

\subsubsection{10. Effective comparison group. {[}5
points{]}}\label{effective-comparison-group.-5-points}

The authors argue that matching can fail even when the observed
covariates are very rich because people who are reachable by phone may
differ from those who are not in ways that are not measured in the data.

Using the logic of the matching estimator, explain what counterfactual
quantity the matched controls are intended to provide for the treated
group.

Then explain why differences in latent turnout propensity between
reachable and unreachable individuals would bias the matching estimate.

\begin{center}\rule{0.5\linewidth}{0.5pt}\end{center}

\section{Part IV. Reveal and reproduce the experimental
benchmark}\label{part-iv.-reveal-and-reproduce-the-experimental-benchmark}

You may now use the randomized treatment-assignment variable.

\subsubsection{11. Experimental benchmark. {[}10
points{]}}\label{experimental-benchmark.-10-points}

Estimate the experimental effect of the GOTV intervention using
randomized assignment.

Report:

\begin{itemize}
\tightlist
\item
  the intent-to-treat effect (ITT);
\item
  the first-stage effect of assignment on phone contact; and
\item
  the treatment-on-treated estimate using assignment as an instrumental
  variable.
\end{itemize}

Explain why randomized assignment provides identification of the
relevant causal effect despite imperfect compliance with treatment
assignment.

Compare the experimental estimate with the observational matching
estimate.

\subsubsection{12. The placebo benchmark. {[}5
points{]}}\label{the-placebo-benchmark.-5-points}

Repeat the experimental analysis for the buckle-up condition.

Explain what the experimental estimate tells us about the plausibility
of the matching estimate for the placebo treatment.

\begin{center}\rule{0.5\linewidth}{0.5pt}\end{center}

\section{Part V. Explain the failure of
matching}\label{part-v.-explain-the-failure-of-matching}

\subsubsection{13. Decompose the source of bias. {[}10
points{]}}\label{decompose-the-source-of-bias.-10-points}

The paper derives an expression showing that matching bias depends on
the difference between the untreated turnout rate of reachable people
and the untreated turnout rate of people who are not reachable.

Using the notation in the paper, explain the intuition behind this
expression without simply restating the algebra.

In particular, explain:

\begin{itemize}
\tightlist
\item
  why reachable untreated individuals provide useful information about
  the counterfactual turnout of treated individuals;
\item
  why unreachable untreated individuals can contaminate the comparison;
  and
\item
  why the fraction of the control group that is reachable matters for
  the size of the bias.
\end{itemize}

\subsubsection{14. Sensitivity analysis. {[}10
points{]}}\label{sensitivity-analysis.-10-points}

Reproduce the paper's sensitivity analysis regarding the amount of
hidden bias necessary to explain the observed matching estimate.

Present the analysis graphically or in a clearly labeled table.

Then interpret the result in substantive terms. How large would the
difference in untreated turnout between reachable and unreachable
individuals have to be to account for the observed matching estimate?

Explain how this analysis changes your assessment of the observational
evidence.

\begin{center}\rule{0.5\linewidth}{0.5pt}\end{center}

\section{Part VI. Causal audit}\label{part-vi.-causal-audit}

\subsubsection{15. Referee report. {[}15
points{]}}\label{referee-report.-15-points}

Write a short referee-style assessment of the following claim:

\begin{quote}
``Because the treated and comparison groups are exactly balanced on a
rich set of pre-treatment covariates, the matching estimate provides a
credible estimate of the causal effect of GOTV phone calls on turnout.''
\end{quote}

Your assessment should distinguish among:

\begin{enumerate}
\def\labelenumi{\arabic{enumi}.}
\tightlist
\item
  evidence about observed covariate balance;
\item
  evidence about overlap and the construction of the comparison group;
\item
  placebo and pre-treatment outcome diagnostics;
\item
  sensitivity to hidden bias; and
\item
  the conditional ignorability assumption itself.
\end{enumerate}

State the strongest causal conclusion you believe could have been
defended \textbf{without access to the randomized benchmark}.

Then explain how the experimental benchmark changes your assessment.

Do not simply say that ``matching is biased.'' Your task is to identify
the specific identifying assumption that fails and explain how the
empirical evidence reveals the failure.

\begin{center}\rule{0.5\linewidth}{0.5pt}\end{center}

\section{Optional extension: Alternative
estimators}\label{optional-extension-alternative-estimators}

This section is optional and will not count toward the core grade.

\subsubsection{A. Regression adjustment}\label{a.-regression-adjustment}

Estimate the GOTV effect using regression adjustment with the same
observed covariates used for matching. Compare the regression estimate
with the matching estimate.

\subsubsection{B. Weighting}\label{b.-weighting}

Estimate the effect using a covariate-balancing or overlap-weighting
estimator. Examine covariate balance before and after weighting and
compare the resulting estimate with the matching and regression
estimates.

\subsubsection{C. Estimator agreement versus
identification}\label{c.-estimator-agreement-versus-identification}

Suppose that regression, matching, and weighting all produced nearly
identical estimates. Would that provide evidence that conditional
ignorability is true? Explain why or why not.

This question is intended to reinforce a central principle of the
course:

\begin{quote}
\textbf{Robustness across estimators does not establish robustness of
identification assumptions.}
\end{quote}

\begin{center}\rule{0.5\linewidth}{0.5pt}\end{center}

\section{References}\label{references}

Arceneaux, Kevin, Alan S. Gerber, and Donald P. Green. 2010. ``A
Cautionary Note on the Use of Matching to Estimate Causal Effects: An
Empirical Example Comparing Matching Estimates to an Experimental
Benchmark.'' \emph{Sociological Methods \& Research}.

Course readings on causal identification, conditioning, matching, and
weighting may be cited where relevant.

\end{document}
